%============================================================
%  PLAN DE TRABAJO PROFESIONAL DE INGENIERÍA ELECTRÓNICA
%  Radar FMCW de 1 a 2 GHz (GPR) — Martín y Santiago
%
%  Esqueleto para escribir a mano. Cada sección tiene un
%  \cartel{tipo de texto}{de qué tiene que tratar}: un recuadro rojo
%  que se reemplaza por el texto. Los datos sueltos que faltan están
%  marcados con \pendiente{...}. Antes de presentarlo: buscar
%  "\cartel" y "\pendiente" y no tiene que quedar ninguno.
%
%  Compilar con compilar.bat (pdflatex + biber + pdflatex x2).
%============================================================

\documentclass[11pt, a4paper, oneside]{article}

%------------------------------------------------------------
% PAQUETES
%------------------------------------------------------------
\usepackage[utf8]{inputenc}
\usepackage[T1]{fontenc}
\usepackage[spanish, es-tabla, es-noshorthands]{babel}
\usepackage{lmodern}
\usepackage{microtype}

% Márgenes amplios y aire entre líneas
\usepackage[
  top=3.0cm, bottom=3.0cm,
  left=3.0cm, right=2.5cm,
  headheight=14pt, headsep=1.0cm, footskip=1.4cm
]{geometry}
\linespread{1.2}
\setlength{\parindent}{0pt}
\setlength{\parskip}{0.7em}

\usepackage{amsmath}
\usepackage[table]{xcolor}
\usepackage{graphicx}
\usepackage{adjustbox}
\usepackage{booktabs}
\usepackage{tabularx}
\usepackage{array}
\usepackage{caption}
\captionsetup{font=small, labelfont=bf, margin=0.5cm}
\usepackage{enumitem}
\setlist{itemsep=0.35em, topsep=0.4em, parsep=0pt}
\setlist[itemize]{leftmargin=1.6em}
\setlist[enumerate]{leftmargin=1.8em}
\usepackage{tikz}
\usetikzlibrary{arrows.meta, positioning, calc, babel}

% Celeste institucional de FIUBA
\definecolor{fiuba}{RGB}{0,162,227}

% Títulos de sección: número con punto, en celeste, con aire
\usepackage{titlesec}
\titleformat{\section}{\Large\bfseries}{\textcolor{fiuba}{\thesection.}}{0.6em}{}
\titleformat{\subsection}{\large\bfseries}{\textcolor{fiuba}{\thesubsection.}}{0.6em}{}
\titleformat{\subsubsection}{\normalsize\bfseries}{\thesubsubsection.}{0.6em}{}
\titlespacing*{\section}{0pt}{2.4em}{1.0em}
\titlespacing*{\subsection}{0pt}{1.8em}{0.6em}
\titlespacing*{\subsubsection}{0pt}{1.4em}{0.4em}

% Encabezado con el nombre de la sección y pie con el número de página
\usepackage{fancyhdr}
\pagestyle{fancy}
\fancyhf{}
\fancyhead[L]{\small\textcolor{black!60}{Plan de Trabajo Profesional}}
\fancyhead[R]{\small\nouppercase{\seccionactual}}
\fancyfoot[C]{\small\thepage}
\renewcommand{\headrulewidth}{0.4pt}
\renewcommand{\sectionmark}[1]{\markboth{\thesection.\ #1}{}}
% La sección vigente al principio de la página; si no hay ninguna
% (la primera), la que arranca en ella
\makeatletter
\newcommand{\seccionactual}{%
  \protected@edef\@seccionarriba{\TopMark{2e-left}}%
  \ifx\@seccionarriba\@empty\FirstMark{2e-left}\else\@seccionarriba\fi}
\makeatother

% Carátula: solo la línea del pie con el tipo de documento
\fancypagestyle{caratula}{%
  \fancyhf{}%
  \renewcommand{\headrulewidth}{0pt}%
  \renewcommand{\footrulewidth}{0.4pt}%
  \fancyfoot[L]{\small Plan de Trabajo Profesional de Ingeniería Electrónica}%
}

\usepackage[
  colorlinks=true,
  linkcolor=black,
  citecolor=fiuba!70!black,
  urlcolor=fiuba!70!black
]{hyperref}

\usepackage{csquotes}
\usepackage[backend=biber, style=ieee, sorting=none]{biblatex}
\addbibresource{referencias_plan.bib}
\AtBeginBibliography{\small\linespread{1.1}\selectfont\setlength{\parskip}{0pt}}
\setlength{\bibitemsep}{0.45em}

%------------------------------------------------------------
% COMANDOS PROPIOS
%------------------------------------------------------------
% Lo que falta completar: en rojo, para que no se escape
\newcommand{\pendiente}[1]{\textcolor{red!75!black}{[#1]}}
% Recuadro rojo con el tipo de texto que va en ese lugar y de qué trata
\usepackage[most]{tcolorbox}
\newtcolorbox{cartelbox}[1]{
  enhanced, breakable,
  colback=red!3, colframe=red!65!black, colbacktitle=red!65!black,
  coltitle=white, fonttitle=\small\bfseries,
  title={A completar: #1},
  fontupper=\small\color{red!40!black},
  boxrule=0.6pt, arc=2pt,
  left=8pt, right=8pt, top=5pt, bottom=6pt,
  before skip=0.5em, after skip=0.9em
}
\newcommand{\cartel}[2]{\begin{cartelbox}{#1}#2\end{cartelbox}}
% Sigla con su significado al pie, la primera vez que aparece
\newcommand{\sigla}[2]{#1\footnote{#2.}}

%------------------------------------------------------------
% DATOS DEL TRABAJO
%------------------------------------------------------------
\newcommand{\Titulo}{Desarrollo y caracterización de un radar FMCW
de 1 a 2~GHz para la detección de objetos ocultos}

%============================================================
\begin{document}

%============================================================
%  CARÁTULA
%============================================================
\begin{titlepage}
  \thispagestyle{caratula}
  \centering
  \vspace*{-0.5cm}

  \adjincludegraphics[width=0.62\textwidth, trim={{.105\width} {.33\height} {.10\width} {.34\height}}, clip]{figuras/logo_fiuba.png}

  \vspace{0.9cm}
  {\large\textsc{Universidad de Buenos Aires}}\\[0.15cm]
  {\large\textsc{Facultad de Ingeniería}}\\[0.15cm]
  {\small Año 2026}

  \vspace{1.6cm}
  {\large\bfseries PLAN DE TRABAJO PROFESIONAL DE INGENIERÍA ELECTRÓNICA}

  \vspace{2.2cm}
  {\Large\linespread{1.35}\selectfont \MakeUppercase{\Titulo}\par}

  \vfill

  {\small\bfseries INTEGRANTES:}\\[0.35cm]
  \pendiente{Apellido}, Martín\\
  Padrón: \pendiente{padrón}\\
  \href{mailto:maledesma@fi.uba.ar}{maledesma@fi.uba.ar}\\[0.45cm]
  Mogica, Santiago\\
  Padrón: \pendiente{padrón}\\
  \pendiente{mail @fi.uba.ar}

  \vspace{0.9cm}
  {\small\bfseries DIRECTOR:}\\[0.35cm]
  Dr. Ing. Fano, Walter Gustavo\\
  \pendiente{mail}

  % Si hay codirector, descomentar:
  % \vspace{0.9cm}
  % {\small\bfseries CODIRECTOR:}\\[0.35cm]
  % Apellido, Nombre\\
  % mail

  \vspace{1.2cm}
\end{titlepage}

%============================================================
%  ÍNDICE
%============================================================
\setcounter{page}{2}
{\setlength{\parskip}{0.25em}\tableofcontents}
\markboth{}{}
\clearpage

%============================================================
\section{Objeto y área del Trabajo Profesional}
%============================================================

\cartel{Presentación del trabajo}{Un párrafo: qué se va a hacer, sobre
qué sistema y con qué fin. Mencionar que se parte de un prototipo
recibido en estado parcial. Las siglas se definen al pie la primera vez
que aparecen (FMCW, GPR, VCO\ldots).}

\subsection{Objetivo general}

\cartel{Objetivo general}{Una o dos oraciones. Mantenerlo general: la
aplicación final (suelo, paredes u otra) se define con los ensayos.}

\subsection{Objetivos particulares}

\cartel{Objetivos particulares}{Lista en viñetas, de 5 a 8 ítems, cada
uno verificable. Por ejemplo, uno por subsistema (barrido, cadena de RF,
digitalización, software, simulación, integración) y uno de evaluación
experimental.}

\subsection{Área del Trabajo Profesional}

\cartel{Áreas profesionales}{Nombrar las áreas de la carrera
involucradas (por ejemplo telecomunicaciones, sistemas embebidos,
procesamiento de señales) y aclarar que están dentro de los alcances
del título de Ingeniero Electrónico.}

%============================================================
\section{Introducción y antecedentes}
%============================================================

\subsection{Motivación}

\cartel{Motivación}{Breve, uno o dos párrafos: para qué sirve un GPR,
por qué conviene uno de bajo costo y por qué la arquitectura FMCW lo
hace viable. Con citas.}

\subsection{Antecedentes}

\subsubsection{Contexto del proyecto}

\cartel{Contexto institucional y trabajo previo}{Quién empezó el
prototipo y en qué laboratorio o grupo, qué componentes se heredaron y
en qué estado se recibió (sin digitalización, barrido con generador de
laboratorio, sin estimación de distancia).}

\subsubsection{Marco normativo}

\cartel{Normas aplicables}{Normas sobre radares de penetración
terrestre y de pared, uso del espectro en la Argentina y exposición
humana a radiofrecuencia. Aclarar que no se busca certificar el equipo
y que las normas se usan para definir condiciones seguras de ensayo.}

%============================================================
\section{Desarrollo previsto del Trabajo Profesional}
%============================================================

\subsection{Teoría, enfoque y métodos a utilizar}

\subsubsection{Principio de funcionamiento}

\cartel{Introducción teórica}{Principio del radar FMCW en uno o dos
párrafos: barrido lineal, retardo del eco y frecuencia de batido.
Incluir la ecuación de la frecuencia de batido y la de la resolución en
distancia, con los valores del sistema (ancho de banda de 1~GHz).}

\subsubsection{Enfoque}

\cartel{Enfoque y metodología}{Cómo se va a encarar el desarrollo (por
subsistemas, con integración incremental). Dejar claro qué se hereda,
qué son módulos comerciales y qué es desarrollo propio. Referenciar la
Figura~\ref{fig:bloques}.}

\begin{figure}[htbp]
  \centering
  \begin{tikzpicture}[
      font=\footnotesize,
      >={Stealth[length=2mm]},
      line width=0.6pt,
      blk/.style={draw, rounded corners=2pt, minimum width=1.6cm,
                  minimum height=1.05cm, inner xsep=2pt, align=center,
                  line width=0.6pt},
      her/.style={blk, draw=black!45, fill=black!8},
      pro/.style={blk, draw=fiuba, fill=fiuba!12},
      sub/.style={font=\scriptsize, text=black!70},
      lab/.style={font=\scriptsize, midway, above, inner sep=1.5pt}
    ]
    \def\dx{2.3}
    \def\yb{-2.3}
    % Transmisión
    \node[pro] (gen) at (0,0)      {Generador\\[-2pt]{\scriptsize triangular}};
    \node[her] (amp) at (\dx,0)    {AMP\\[-2pt]{\scriptsize $\times 5$}};
    \node[her] (vco) at (2*\dx,0)  {VCO\\[-2pt]{\scriptsize 1--2 GHz}};
    \node[her] (att) at (3*\dx,0)  {Atenuador\\[-2pt]{\scriptsize $-3$ dB}};
    \node[her] (spl) at (4*\dx,0)  {Divisor de\\[-2pt]potencia};
    \node[her] (tx)  at (5*\dx,0)  {Antena\\[-2pt]TX};
    % Recepción
    \node[pro] (esp) at (0,\yb)     {ESP32-C3\\[-2pt]{\scriptsize firmware}};
    \node[pro] (pcm) at (\dx,\yb)   {PCM1808\\[-2pt]{\scriptsize ADC 24 bit}};
    \node[her] (lpf) at (2*\dx,\yb) {Filtro\\[-2pt]{\scriptsize pasabajos}};
    \node[her] (mix) at (3*\dx,\yb) {Mezclador};
    \node[her] (lna) at (4*\dx,\yb) {LNA};
    \node[her] (rx)  at (5*\dx,\yb) {Antena\\[-2pt]RX};
    % Medio y PC
    \node[blk, dashed, draw=black!60, minimum width=1.4cm] (med) at (5*\dx+1.95,\yb/2) {Medio\\[-2pt]{\scriptsize bajo estudio}};
    \node[pro] (pc) at (0,\yb-1.6) {PC\\[-2pt]{\scriptsize procesamiento}};
    % Conexiones de transmisión
    \draw[->] (gen) -- (amp);
    \draw[->] (amp) -- (vco);
    \draw[->] (vco) -- (att);
    \draw[->] (att) -- (spl);
    \draw[->] (spl) -- (tx);
    \draw[->, dashed] (tx.east) -| (med.north);
    \draw[->, dashed] (med.south) |- (rx.east);
    % Conexiones de recepción
    \draw[->] (rx) -- (lna);
    \draw[->] (lna) -- node[lab]{RF} (mix);
    \draw[->] (spl.south) -- ++(0,-0.62) -| node[pos=0.25, above, font=\scriptsize]{LO} (mix.north);
    \draw[->] (mix) -- node[lab]{$f_b$} (lpf);
    \draw[->] (lpf) -- (pcm);
    \draw[->] (pcm) -- node[lab]{I$^2$S} (esp);
    \draw[->] (gen) -- node[midway, left, font=\scriptsize]{sync} (esp);
    \draw[->] (esp) -- node[midway, right, font=\scriptsize]{USB} (pc);
    % Leyenda
    \node[her, minimum width=0.55cm, minimum height=0.35cm] (lh) at (2.6,\yb-1.6) {};
    \node[right=0.15cm of lh, font=\footnotesize] {Heredado del prototipo previo};
    \node[pro, minimum width=0.55cm, minimum height=0.35cm] (lp) at (8.4,\yb-1.6) {};
    \node[right=0.15cm of lp, font=\footnotesize, align=left] {Desarrollo propio\\[-2pt]{\scriptsize (incluye módulos comerciales integrados)}};
  \end{tikzpicture}
  \caption{Diagrama de bloques del radar. En gris, los componentes
  recibidos del prototipo previo; en celeste, los desarrollados o
  integrados en este trabajo.}
  \label{fig:bloques}
\end{figure}

\cartel{Revisar el diagrama}{Confirmar los bloques y qué es heredado
(gris) y qué es propio (celeste).}

\subsubsection{Generación del barrido}

\cartel{Hardware diseñado: generador del barrido}{Qué hace el bloque,
cómo se implementa (generador de triangular analógico) y qué problema
resuelve (alinealidad del VCO). Cerrar con las tareas previstas en
viñetas.}

\subsubsection{Cadena de radiofrecuencia y antenas}

\cartel{Hardware heredado: cadena de RF y antenas}{Descripción breve de
la cadena recibida y aclaración de que no se rediseña. Tareas previstas
en viñetas: caracterización de cables con VNA, potencias, filtro y,
si corresponde, antenas.}

\subsubsection{Recepción y digitalización}

\cartel{Hardware y firmware diseñados: digitalización}{ADC y
microcontrolador elegidos, cómo se conectan y cómo llegan los datos a
la PC. Tareas previstas en viñetas.}

\subsubsection{Software desarrollado}

\cartel{Software diseñado}{Qué hace el software propio: adquisición,
corrección de la alinealidad, estimación de distancia, calibración y
visualización en tiempo real. Tareas previstas en viñetas.}

\subsubsection{Herramientas utilizadas}

\cartel{Software y herramientas utilizados}{Herramientas de terceros y
para qué se usa cada una: por ejemplo Python, LTspice, KiCad, MEEP, CST
y el entorno del microcontrolador, además del instrumental de
laboratorio (osciloscopio, analizador de espectro, VNA).}

\subsubsection{Simulación electromagnética}

\cartel{Simulación electromagnética}{Qué se simula (escenarios de
ensayo, antenas, blancos), con qué método y herramienta, y para qué se
usa (contrastar con las mediciones).}

\subsubsection{Verificación}

\cartel{Criterio de verificación}{Verificación progresiva: subsistema,
integración y sistema completo. Cada resultado se contrasta por tres
vías: cálculo analítico, simulación y medición. Ensayos de laboratorio
y de campo.}

\subsection{Estudios conexos}

\cartel{Estudios conexos}{Materias de la carrera relacionadas con el
trabajo, agrupadas, con una línea sobre qué aporta cada grupo.
Opcional: formación extra (cursos, trabajo).}

\subsection{Alcance proyectado para el Trabajo Profesional}

\cartel{Resultados esperados}{Una oración de introducción (el alcance es
general y la aplicación se define con los ensayos) y la lista de
resultados concretos en viñetas, con el nombre de cada resultado en
negrita.}

\cartel{Aplicación a evaluar}{Orden de prioridad: 1) detección de
objetos enterrados; 2) a través de paredes; 3) otro uso, por ejemplo
SAR. Cada alternativa se aborda solo si la anterior no resulta viable.
Indicar en qué revisión se decide.}

\cartel{Qué no incluye}{Un párrafo con lo que queda fuera: rediseño de
los componentes heredados, algoritmos avanzados de imagen, una campaña
exhaustiva de suelos y la certificación del equipo.}

\subsection{Plan de trabajo}
\label{sec:plan}

\cartel{Presentación del plan}{Duración (43 semanas, unos 10 meses,
desde el 3 de agosto de 2026), esfuerzo total del equipo (unas
1200~h), revisiones formales al cierre de cada etapa y redacción
continua desde el principio. Referenciar las
Tablas~\ref{tab:gantt} y~\ref{tab:etapas}.}

\cartel{Revisar las tablas}{Las etapas, las semanas, las horas y las
revisiones de las dos tablas son una propuesta inicial.}

\newcommand{\hc}[1][]{\cellcolor{fiuba}\textcolor{white}{\scriptsize\bfseries #1}}
\newcommand{\pc}[1][]{\cellcolor{fiuba!25}{\scriptsize\bfseries #1}}

\begin{table}[htbp]
  \centering
  \small
  \setlength{\tabcolsep}{0pt}
  \renewcommand{\arraystretch}{1.5}
  \newcolumntype{M}{>{\centering\arraybackslash}p{0.95cm}}
  \begin{tabular}{@{} >{\raggedright\arraybackslash}p{5.0cm} | M M M M M | M M M M M @{}}
    & \multicolumn{5}{c|}{\textbf{2026}} & \multicolumn{5}{c}{\textbf{2027}} \\
    \textbf{Etapa} & Ago & Sep & Oct & Nov & Dic & Ene & Feb & Mar & Abr & May \\
    \midrule
    E1 -- Banco de adquisición    & \hc[R1] & & & & & & & & & \\
    E2 -- Cadena heredada         & \hc & \hc[R2] & & & & & & & & \\
    E3 -- Mediciones y simulación & & \hc & \hc[R3] & & & & & & & \\
    E4 -- Generador del barrido   & & \hc & \hc & \pc[R4] & & & & & & \\
    E5 -- Integración             & & & & \pc & \pc[R5] & & & & & \\
    E6 -- Ensayos en suelo        & & & & & \pc & \pc & \pc[R6] & & & \\
    E7 -- Aplicación elegida      & & & & & & & \pc & \pc & \pc[R7] & \\
    E8 -- Documentación final     & & & & & & & & & \pc & \pc[R8] \\
    \midrule
    Redacción del informe         & \hc & \hc & \hc & \pc & \pc & \pc & \pc & \pc & \pc & \pc \\
    \bottomrule
  \end{tabular}

  \vspace{0.6em}
  {\footnotesize
   \tikz\fill[fiuba] (0,0) rectangle (0.35,0.22); \ cumplido
   \qquad
   \tikz\fill[fiuba!25] (0,0) rectangle (0.35,0.22); \ previsto
   \qquad
   R$n$: revisión de cierre de la etapa}
  \caption{Cronograma mensual de las etapas.}
  \label{tab:gantt}
\end{table}

\cartel{Descripción de las etapas}{Lista numerada de E1 a E8, con dos o
tres líneas por etapa: qué se hace y con qué revisión cierra.}

\begin{table}[htbp]
  \centering
  \small
  \renewcommand{\arraystretch}{1.35}
  \begin{tabularx}{\textwidth}{@{} l c X c r @{}}
    \toprule
    \textbf{Etapa} & \textbf{Semanas} & \textbf{Revisión de cierre}
      & \textbf{Estado} & \textbf{Esfuerzo} \\
    \midrule
    E1 -- Banco de adquisición        & 1--4   & R1: banco validado                   & Cumplida & 120 h \\
    E2 -- Cadena heredada             & 3--7   & R2: cadena caracterizada             & Cumplida & 150 h \\
    E3 -- Mediciones y simulación     & 5--10  & R3: distancia contra simulación      & Cumplida & 160 h \\
    E4 -- Generador del barrido       & 8--17  & R4: generador verificado             & En curso & 110 h \\
    E5 -- Integración                 & 14--22 & R5: prototipo integrado              & Prevista & 160 h \\
    E6 -- Ensayos en suelo            & 21--28 & R6: definición de la aplicación      & Prevista & 170 h \\
    E7 -- Aplicación elegida          & 29--37 & R7: resultados experimentales        & Prevista & 180 h \\
    E8 -- Documentación final         & 36--43 & R8: entrega del informe              & Prevista & 150 h \\
    \midrule
    \textbf{Total} & \textbf{43} & & & \textbf{1200 h} \\
    \bottomrule
  \end{tabularx}
  \caption{Etapas, revisiones y esfuerzo estimado del equipo. La
  semana 1 comienza el 3 de agosto de 2026.}
  \label{tab:etapas}
\end{table}

%============================================================
\section{Referencias}
%============================================================

\cartel{Referencias}{Se listan solas las entradas de
\texttt{referencias\_plan.bib} citadas en el texto con
\texttt{\textbackslash cite}. El archivo ya trae algunas sugeridas;
borrar las que no se usen.}

\printbibliography[heading=none]

%============================================================
%  FIRMAS
%============================================================
\clearpage
\thispagestyle{empty}
\vspace*{\fill}

\begin{flushright}
  Buenos Aires, \pendiente{fecha de presentación}
\end{flushright}

\vspace{3.5cm}

\begin{center}
  \begin{tabular}{@{} c @{\hspace{2.5cm}} c @{}}
    \rule{6cm}{0.4pt} & \rule{6cm}{0.4pt} \\[0.2cm]
    \pendiente{Apellido}, Martín & Mogica, Santiago \\
    \small Integrante & \small Integrante \\
  \end{tabular}

  \vspace{3.5cm}

  \begin{tabular}{@{} c @{}}
    \rule{6cm}{0.4pt} \\[0.2cm]
    Fano, Walter Gustavo \\
    \small Director \\
  \end{tabular}
\end{center}

\vspace*{\fill}

\end{document}
